\section{国内外研究现状}
\label{intro:related}

无人机协同通信受到拓扑、能量、消息内容和反馈状态的共同影响。现有研究分别从路由、在线资源调度、协同感知、语义表示和可靠传输等角度展开。其中，路由研究关注可用路径与传输代价，协同感知研究关注信息交换对判断结果的作用，可靠传输研究则进一步处理接收结果不确定时的后续动作。本节沿这些方向梳理已有方法，分析其解决的问题及适用条件。

\subsection{无人机自组网路由与能量控制}

飞行自组网中的节点既执行感知任务，也承担数据中继。Bekmezci等指出，三维空间中的快速运动、频繁变化的连接关系与受限机载资源，是飞行自组网区别于传统移动网络的基本因素\cite{bekmezci2013fanet}。Zeng等分析了无人机通信的体系结构和信道特征\cite{zeng2016uav}；Mozaffari等进一步从空中基站与蜂窝终端两类角色讨论部署、轨迹和资源分配\cite{mozaffari2019tutorial}。无人机的灵活部署能够改善覆盖，但同时要求协议适应运动和能量变化。

传统路由主要依赖拓扑或地理位置确定下一跳。AODV按需发现路由，利用目的序列号判断路径的新旧程度，减少无业务时的维护开销\cite{aodv}；OLSR主动维护拓扑，通过多点中继减少控制消息的重复转发\cite{olsr}。两者在路径发现时机与维护开销之间采用不同取舍，节点快速移动时，已有路径都需要随连接关系更新。GPSR根据邻居与目的节点的位置贪婪转发，在局部无法前进时切换到周边转发\cite{gpsr}。地理路由利用局部信息降低了对完整拓扑的依赖，但位置进度不能直接表示传输成功率与实际能耗。这些方法具有明确的转发规则，却可能使流量长期集中于位置有利的少数节点，路径可用性并不能保证能量均衡。

针对节点高速移动，后续方法引入链路稳定性、运动趋势和剩余能量。Lakew等归纳了拓扑变化、维护开销与网络寿命之间的制约关系\cite{intro4}。QGeo利用Q学习修正地理路由中的下一跳价值\cite{qgeo}，QTAR进一步将拓扑状态纳入学习过程\cite{qtar}，从而根据交互结果调整转发选择。与固定规则相比，学习方法可以把历史通信结果用于候选评价，但其优化方向仍取决于状态和回报中是否包含相应的资源因素。Li等提出节能快速路由方法，结合电池、信道和邻居状态联合选择下一跳与发射功率\cite{drlfr}，使单次传输的收益与实际代价共同参与决策。

路由与功率存在耦合：提高功率可能改善成功率，也会增加能耗与干扰；降低功率虽然节省单次支出，却可能产生更多重传。连续任务还要求考虑历史消耗，避免当前收益较高的路径因反复使用而提前失效。本文第一项研究由此引入长期能量反馈，在局部路由中协调转发效率与节点持续服务能力。

\subsection{基于队列稳定性的在线资源调度}

业务到达和无线信道难以事先准确预测，依赖完整未来状态的资源分配不易用于持续运行系统。Tassiulas和Ephremides建立了受约束排队系统的稳定性分析基础，利用队列积压表示尚未满足的服务需求\cite{tassiulas1992stability}。队列随业务到达增加、随实际服务减少，因而能够反映当前负载与服务能力之间的差额。Neely系统研究了李雅普诺夫漂移及漂移加惩罚方法，将长期随机约束转化为依赖当前队列状态的逐时隙决策\cite{neely2010stochastic}。对于平均能量等约束，可构造虚拟队列累计支出偏差，使历史资源压力持续影响后续动作。虚拟队列不代表实际缓存，而是把多次决策造成的资源偏差保留在状态中；当消耗持续超过目标时，后续控制会提高对能量代价的重视程度。

反压路由根据邻居之间的队列差调节流量，能够响应负载变化，但可能因绕行增加时延。Ying等结合最短路径和反压路由，以缩短转发路径\cite{related2}；Yu和Neely研究联合速率控制与路由，分析效用最优性差距与有限队列长度的关系\cite{algo1}。这些工作说明，队列稳定性需要与业务时延和路径代价共同考虑。在飞行自组网中，负载较低的邻居还可能具有较差链路或较少剩余能量，仅按队列差分流不能覆盖全部转发条件，需要与可行链路集合及通信代价结合。

在线队列方法也被用于无人机资源管理。Lin等面向太阳能无人机边缘计算，联合优化资源分配和三维轨迹，在队列与能量约束下组织感知、卸载和通信\cite{intro2}。不过，能量采集、轨迹可控性以及全局状态的可获得程度，都会改变时隙子问题。面向可充电系统的平均收支约束，与固定电池下任务结束前的可用能量约束也不能直接等同。本文主要利用虚拟队列保留资源消耗的时间关联：LEADR在连续时隙中调节路由，C-BFC在有限任务序列中调节交付代价，两者的约束性质分别在对应模型下讨论。

\subsection{多智能体选择性通信与协同感知}

多智能体通信用于弥补单个节点观察范围的不足。RIAL和DIAL将消息与任务动作共同学习，DIAL在训练时允许梯度通过通信连接传播\cite{rial-dial}；CommNet使用连续消息连接智能体隐状态\cite{commnet}。这类方法能够从任务反馈中学习通信内容，但频繁的全体交换会使通信量随节点规模增长。

后续研究因而显式选择通信对象和时机。IC3Net以门控决定是否发送\cite{ic3net}，SchedNet在共享介质和并发发送限制下学习广播调度\cite{schednet}，TarMAC则学习向哪些接收者传递消息\cite{tarmac}。三者分别对应是否通信、由谁通信和向谁通信，使资源分配成为可学习的动作。其共同特点是把通信选择嵌入任务策略，依据协作收益训练消息或门控。当消息具有不同长度、可靠性和表达粒度时，是否发送还需与发送成本一同评价，仅限制同时发送的节点数量不足以描述实际资源支出。

在协同感知中，观察差异通常具有空间含义。V2VNet交换压缩中间特征，用于联合感知与预测\cite{v2vnet}；DiscoNet借助知识蒸馏学习协作图，对不同区域的邻居特征赋予不同权重\cite{disconet}。中间特征可减少原始观测传输，但仍需处理视野重叠造成的冗余。Who2com通过学习握手匹配请求方和信息提供方\cite{who2com}，When2com共同学习协作关系与通信时机\cite{when2com}；Where2comm利用置信图筛选稀疏空间特征\cite{where2comm}，How2comm进一步结合空间与通道选择、时延补偿和时空融合\cite{how2comm}。前两者侧重选择协作节点及通信关系，后两者将选择范围细化到空间区域和特征通道，能够更有针对性地保留互补信息。它们也说明，信息价值与接收者的已有观察密切相关，单独提高发送方的局部置信度并不等于提高协作收益。

按空间位置筛选信息适合具有区域对应关系的特征融合。通信对象若转为任务标签和摘要，其作用还取决于表达中保留了哪些判断依据，同样长度的两条消息也可能覆盖不同语义。因而，空间稀疏性能够指导减少冗余，但还需要结合任务输出与消息粒度评价传输收益。

消息选择还具有顺序依赖：一次融合可能改变后续证据的作用，仅依赖静态置信度难以判断何时停止。空间上的互补关系能够帮助发现缺失区域，但不能完全说明一条新消息能否修正当前判断，尤其当候选包含不同粒度的标签和摘要时，重复语义与不同来源的一致支持还需要区别对待。图注意力网络通过关系权重表示候选之间的联系\cite{gat}，强化学习则提供了根据长期回报学习选择序列的工具\cite{sutton2018rl}。DQN估计离散动作价值，Double DQN分离选择与评估以缓解过高估计\cite{mnih2015dqn,vanhasselt2016double}；PPO以裁剪目标抑制过大的策略更新，结合广义优势估计改善训练稳定性\cite{ppo,gae}。这些工具提供表示和求解方式，仍需根据具体通信任务确定状态、动作与回报。本文第二项研究结合图结构和通信历史，在候选、中心需求与剩余预算的共同状态下学习证据选择。

\subsection{任务语义通信与视觉语言证据}

经典信息论用熵、信道容量和率失真描述有限资源下的信息恢复能力\cite{shannon1948,cover2006information}。协同判断通常不要求重建全部观测，而需要保留支持决策的信息。信息瓶颈通过约束中间表示的冗余并保留目标相关性，为任务压缩提供了分析思路\cite{tishby2000ib}。

深度联合信源信道编码把编码器、信道与解码器联合训练。Bourtsoulatze等学习图像到信道符号的映射，以改善受限带宽或信噪比下的传输\cite{bourtsoulatze2019deepjscc}；DeepSC利用Transformer学习文本语义表示与传输映射\cite{xie2021deepsc,vaswani2017attention}。Shao等面向边缘推理，以信息瓶颈保留有效特征，并使用可变长度表示适应信道\cite{shao2021task}。G\"und\"uz等强调应结合上下文和接收端功能评价通信效果\cite{semantic-communications-jsac}，使优化目标从恢复精度延伸到任务完成质量。这些方法主要研究给定内容如何形成适合任务和信道的表示；多源协同还需要决定哪些节点的内容应被选入通信过程。表示压缩与来源选择因此是相互联系的两个问题，前者节省单条消息的资源，后者减少缺少新增作用的消息。

视觉语言模型使任务信息能够以显式证据组织。CLIP通过图文对比学习建立共享表示\cite{clip}；BLIP-2以查询Transformer连接冻结视觉编码器和语言模型\cite{blip2}；LLaVA通过视觉指令微调支持围绕图像的回答与解释\cite{llava}。这些能力可用于提取对象、属性和局部判断，减少多节点系统对完整图像交换的依赖。

但同一图像面对不同问题时，所需信息并不相同。Shih等研究视觉问答中的问题相关区域选择\cite{where-to-look-vqa}；VQA v2利用同一问题对应不同答案的图像对减少语言先验偏差\cite{vqa-v2}；GQA以场景图构造组合推理问题\cite{gqa}。这些研究提示，局部观察是否有用需要相对于问题判断，能够生成完整描述并不意味着其中每部分都值得传输。COCO的图像和标注则为局部观测构造提供基础\cite{coco}。当多个节点覆盖同一对象时，还需要保留证据来源与观察范围，才能分析它们是相互补充还是重复表达。本文据此设置任务相关证据层，保留发送者、观察范围及不同消息粒度，在既定语义模型上学习调度。该分工便于分析通信策略本身的作用，也为后续区分新证据和重复发送提供统一对象。

\subsection{不完备反馈下的可靠传输与有限前瞻决策}

消息选定后仍需经过实际链路。自动重传请求根据确认决定是否再次发送，停等、回退和选择重传分别在控制复杂度、缓存与效率之间取舍\cite{lin1984arq}。传统机制以分组接收为主要依据，而语义任务还需要考虑重传能否改善判断：一份正确到达的分组可能作用有限，一份缺失的关键证据却可能决定结论。

已有工作将任务目标引入重传。Sheng等在协同感知中利用重要性图提取特征，结合语义错误检测和HARQ提高衰落信道下的可靠性\cite{sheng2024harq}。Li等研究语义HARQ，在重复语义信号、多轮接收融合和增加语义冗余之间设计不同方案\cite{li2025semharq}。这些方法将重传与感知质量联系起来，使重传内容及多轮融合不再仅围绕比特错误组织。不过，改进一次重传的有效性与判断当前是否需要重传仍有区别。在确认延迟或丢失时，发送端还需估计接收状态，并比较重传已有证据与发送其他证据的机会成本。

部分可观测马尔可夫决策过程以信念分布表示隐藏状态，使动作能够利用不完整反馈和历史记录\cite{kaelbling1998pomdp}。Huang等以信息效用损失刻画收发状态不一致的影响，在卫星融合网络中联合选择采样、编码分组与重传\cite{uoi2024satellite}。这类状态估计方法与语义重传的联系在于：发送收益既由消息内容决定，也由接收方可能已经获得哪些内容决定。对于无人机语义交付，未收到确认可能意味着数据失败，也可能只是确认未到，两种情况需要不同的动作。立即补发可能增加重复传输，持续等待则会占用有限任务窗口，因此反馈不确定性本身需要进入决策状态。

完整信念规划会随候选和反馈分支增加而迅速变大。模型预测控制采用有限时域规划并随新状态滚动更新，为约束下的在线计算提供思路\cite{mayne2000mpc}；离散动作学习也可估计后续收益，SD-SAC针对离散软演员评论家的价值低估和训练不稳定进行改进\cite{zhou2024sdsac}。本文第三项研究采用显式接收信念与有限前瞻，把新发、重传、等待和停止放在同一过程内比较。等待收益由确认能否及时返回、能否改变后续选择共同决定，因此需要随任务状态计算。

\subsection{现有工作的不足与本文研究定位}

在动态路由与资源调度方面，现有方法能够利用位置、链路或历史回报改善传输，但单次代价较低的路径可能被反复使用，关键节点仍会较快消耗电量。队列方法提供了保留长期资源压力的途径，其具体约束与可观测状态仍需结合任务确定。固定通信预算下，如何在局部信息和有限计算量下协调业务积压、链路功耗及节点能量分布，是持续执行任务时需要解决的问题。

在语义通信与协同感知方面，选择来源、空间区域和特征通道能够减少无关信息，但候选作用还会随接收方已有信息变化。任务证据具有不同来源和表达粒度，静态置信度或固定发送数量难以同时反映其可靠性、补充作用及链路成本。因此，需要在任务状态和通信历史下评价候选序列，并根据后续增益决定停止时机。

在语义交付方面，重要性筛选与多轮融合改善了重传内容的利用，但接收状态未知时仍需判断是否继续投入资源。数据失败和确认未到会产生相同的本地表象，却对应不同的重传需求。显式状态估计能够保留这种差异，有限前瞻则需要进一步权衡反馈带来的信息、等待占用的窗口以及后续任务的能量需求。

据此，本文从通信承载、内容选择和交付控制三个层次展开研究，分别处理传输机会的持续性、任务证据的有效性与接收结果的不确定性，在相应模型和实验中分析有限资源下的协同通信取舍。
